\subsection*{如何使用本书}
\addcontentsline{toc}{subsection}{如何使用本书}

本书可以按两种方式阅读。第一种是课程顺序阅读：从推理时计算与 learning to reason 出发，再进入 memory、planning、coding agents、多模态环境、formal mathematics，以及 abstraction 与 security。这样读的好处是可以看到课程如何一步步把 ``reasoning'' 扩展成 ``agent systems''。第二种是按问题路径阅读。如果你更关心 reasoning，可以优先阅读关于 inference-time computation、post-training recipe 和 theorem proving 的章节；如果你更关心 systems，可以优先阅读 coding agents、多模态环境与安全边界；如果你更关心 formal methods，则建议直接把 autoformalization、verification 与 theorem proving 几章连读。

阅读正文时，建议把每一章当作教材章节，而不是视频摘要。章内的写法遵循统一结构：学习目标、背景与问题设置、核心概念、主体讲解、公式或算法解释、readings 连接、失败模式与边界条件、章节小结、复习题与延伸阅读。若某一页 slides 很 dense，本书的预期做法不是把它压缩成一句 takeaway，而是拆开说明图、表、公式与 bullet 各自在论证中的作用。

如果你需要验证某个说法来自哪里，不要只在正文里寻找脚注。每一章背后都保留了 lecture workspace，可回查 \texttt{source\_manifest.json}、\texttt{transcript.jsonl}、\texttt{slides.jsonl}、\texttt{coverage\_units.jsonl}、\texttt{figure\_manifest.json}、\texttt{readings\_integration.md}、\texttt{eval\_report.json} 与 \texttt{lecture\_quality\_report.md}。这些 sidecar 文件记录了源材料、覆盖账本、配图 provenance、reading 融合方式，以及 evaluator/validator 的门禁结果。正文负责可读性，workspace 负责可审计性。

当你看到英文术语时，本书通常会在首次出现时给出中英双语，例如“推理时计算（inference-time computation）”“自动形式化（autoformalization）”“证明搜索（proof search）”。后续章节为了可读性会在中文与英文缩写之间切换，但概念定义应与附录中的 glossary 和 notation table 保持一致。若某一讲必须暂时重载符号或采用讲者自己的术语，相关差异会记录在对应 lecture 的 \texttt{notation\_delta.md} 或 \texttt{glossary\_delta.md} 中。

本书还刻意保留了置信度与遗漏信息。若字幕中某段难以辨认、某篇 reading 只能获取摘要、或课程页提供的信息不足以支持更强断言，相应问题应进入 \texttt{low\_confidence\_spans.jsonl}、\texttt{omission\_log.jsonl} 或课程级 omission appendix，而不是被静默掩盖。阅读时可以把这种设计理解为一条使用约定：本书追求高覆盖和高可读性，但不会用流畅措辞替代证据。

最后，这本书与普通视频笔记最大的差别，在于它不是在完成后才补 provenance，而是先把 provenance、coverage 和 evaluation 写进仓库，再产出成书文本。对研究者和工程实践者来说，这意味着你不仅能读到结论，也能判断这些结论是如何被组织、验证和限制的；而这恰好与本课程讨论 agent systems 的方式保持一致。
